% Part: turing-machines
% Chapter: undecidability
% Section: unsolvability-decision-problem

\documentclass[../../../include/open-logic-section]{subfiles}

\begin{document}

\olfileid[hi]{tur}{und}{dec}
\olsection{निर्णय समस्या}

हम कहते हैं कि प्रथम-कोटि तर्कशास्त्र \emph{निर्णेय} है, तभी और केवल
तभी जब यह निर्धारित करने की कोई प्रभावी विधि हो कि दिया गया
!!{sentence} वैध है या नहीं। वास्तव में ऐसी कोई विधि नहीं है: प्रथम-कोटि
वाक्यों की वैधता का निर्णय करने की समस्या असमाधेय है।

इस महत्त्वपूर्ण निषेधात्मक परिणाम को स्थापित करने के लिए हम सिद्ध करते
हैं कि निर्णय समस्या को ट्यूरिंग मशीन द्वारा हल नहीं किया जा सकता।
अर्थात् हम दिखाते हैं कि ऐसी कोई ट्यूरिंग मशीन नहीं है जो प्रथम-कोटि
!!{sentence} वाली पट्टी पर आरंभ किए जाने पर अंततः रुके और वह !!{sentence}
वैध है या नहीं, इसके अनुसार $1$ अथवा~$0$ निर्गत करे। चर्च--ट्यूरिंग
अभिधारणा के अनुसार प्रत्येक संगणनीय फलन ट्यूरिंग-संगणनीय है। अतः यदि यह
``वैधता फलन'' किसी भी प्रकार प्रभावी रूप से संगणनीय होता, तो वह
ट्यूरिंग-संगणनीय होता। यदि वह ट्यूरिंग-संगणनीय नहीं है, तो वह प्रभावी
रूप से संगणनीय भी नहीं हो सकता।

निर्णय समस्या को असमाधेय सिद्ध करने की हमारी कार्यनीति विराम समस्या को
इसमें अपचयित करना है। इसका अर्थ यह है: हमने सिद्ध किया है कि फलन~$h(e,w)$,
जो निर्गत~$1$ के साथ रुकता है यदि~$e$ द्वारा वर्णित ट्यूरिंग मशीन
आगत~$w$ पर रुकती है और अन्यथा निर्गत~$0$ देता है, ट्यूरिंग-संगणनीय
नहीं है। हम
दिखाएँगे कि यदि प्रथम-कोटि वाक्यों की वैधता का निर्णय करने वाली ट्यूरिंग
मशीन होती, तो~$h$ की संगणना करने वाली ट्यूरिंग मशीन भी होती। चूँकि $h$
की संगणना ट्यूरिंग मशीन द्वारा नहीं की जा सकती, वैधता का निर्णय करने
वाली ट्यूरिंग मशीन भी नहीं हो सकती।

इस कार्यनीति का पहला चरण यह दिखाना है कि प्रत्येक आगत~$w$ और ट्यूरिंग
मशीन~$M$ के लिए हम प्रभावी रूप से !!a{sentence} $!T(M, w)$ का वर्णन
कर सकते हैं, जो~$M$ के निर्देश-समुच्चय और आगत~$w$ को निरूपित करता है,
तथा !!a{sentence}~$!E(M, w)$ का, जो व्यक्त करता है कि “$M$ अंततः रुकती
है,” इस प्रकार कि:
\begin{quote}
  $\Entails !T(M, w) \lif !E(M,w)$ तभी और केवल तभी जब $M$ आगत~$w$ पर
  रुकती है।
\end{quote}
हमारे प्रमाण का अधिकांश भाग इन वाक्यों $!T(M, w)$ और~$!E(M, w)$ का वर्णन
करने तथा यह सत्यापित करने में होगा कि $!T(M, w) \lif !E(M, w)$ तभी और
केवल तभी वैध है जब $M$ आगत~$w$ पर रुकती है।

\end{document}
